This lemma demonstrates the usefulness of monomial ideals. The ability to reduce ideals to finite generators is a step towards closing the gap between the machine and abstract. 

\begin{theorem}
Let $I = \ideal{x^\alpha | \alpha \in A}$ be a monomial ideal. Then $I$ can be finitely written by $I = \ideal{x^{\alpha(1)}, ... , x^{\alpha(s)}}$, where all $\alpha$ lie in $A$.  
\end{theorem}

\begin{proof}
   Consider an inductive approach. Starting off with $n=1$ for our $k[\vectorComponents{x}]$, it is easy to see that the smallest element of $A \subseteq \Z_{\ge 0}$ generates the remaining elements. Whith $n > 1$, we start our induction. Consider the $n$th case. We have the monomial ideal, call it $J$, from the $n-1$ case. This is finite by our induction, so when our $A$ goes from a subset of $\Z_{\ge 0}^{n-1}$ to $\Z_{\ge 0}^{n}$, we can generate slices of these $J$ ideals with the new $x_n = y$ variable. We take some sufficiently large $m \ge 0$ and group the $J_l$ ideals by their corresponding $y$ value up to $m-1$. In other words, let the ideal $J_l$ be generated by $x^\beta$ where $x^\beta y^l \in I$. These $J_l$ slices form a topological map of the ideal $I$. Thus, $I$ is also finite.    
\end{proof}

After a finite ideal is generated, the monomial ideal membership problem is solved. This is since we can take a polynomial and check if the remainder from division by the finite monomial ideal is $0$. If it is, it belongs to the monomial ideal since one of the members will generate it. 